\jlauChapterAfterPart{材料与方法}

\section{示例数据与变量说明}

本示例使用公开可分享的模拟数据结构说明论文写法。数据表包含样方编号、采样日期、冠层覆盖度、归一化植被指数、株高、叶面积指数和环境观测值等字段。真实论文写作时，用户可将本节替换为试验地点、材料来源、田间设计、仪器型号和采样流程。

\section{数据预处理}

预处理流程包括缺失值标记、异常观测复核、时间窗口对齐和变量标准化。对于来自不同设备的观测值，先按设备批次检查分布，再统一进入融合模型，避免设备差异改变结论方向。

\section{融合模型设计}

融合模型由特征筛选、权重估计和稳定性评价三部分组成。示例中用线性加权方式表达基本思路：

\begin{equation}
  F_i=\sum_{j=1}^{p} w_j x_{ij}
\end{equation}

其中，$F_i$ 表示第 $i$ 个样方的融合得分，$x_{ij}$ 表示标准化后的第 $j$ 个特征，$w_j$ 表示特征权重。真实研究可根据课题需要替换为混合效应模型、随机森林、贝叶斯模型或深度学习模型。

\section{图表呈现规范}

图题置于图下方，表题置于表上方，并按章节编号。表格默认采用三线表，以符合规范中对表格编排的要求。

\begin{table}[htbp]
  \centering
  \caption{示例变量说明}
  \begin{tabular}{lll}
    \toprule
    变量 & 含义 & 来源 \\
    \midrule
    CC & 冠层覆盖度 & 图像识别 \\
    NDVI & 归一化植被指数 & 多光谱观测 \\
    LAI & 叶面积指数 & 田间测量 \\
    \bottomrule
  \end{tabular}
\end{table}
